% ============================================================
% Exhibit 4: 七大预期差板块核心指标对比表
% ============================================================
\begin{table}[htbp]
  \centering
  \scriptsize
  \renewcommand{\arraystretch}{1.3}
  \setlength{\tabcolsep}{3.5pt}
  \begin{tabular}{@{}lccccccccccc@{}}
    \toprule
    \textbf{板块名称} & \textbf{景气度} & \textbf{拥挤度} & \textbf{预期差} & \textbf{2026E营收} & \textbf{2026E利润} & \textbf{PE} & \textbf{PE历史} & \textbf{公募持仓} & \textbf{核心催化剂} & \textbf{核心风险} & \textbf{代表标的} \\
    & \textbf{评分} & \textbf{评分} & \textbf{等级} & \textbf{增速(\%)} & \textbf{增速(\%)} & \textbf{(TTM)} & \textbf{分位(\%)} & \textbf{比例(\%)} & & & \\
    \midrule
    \textbf{AI算力} & 88 & 22 & \textcolor{red}{$\bigstar\bigstar\bigstar$} & 42 & 65 & 38 & 18 & 6.8 & 大模型迭代、算力基建 & 算力投资不及预期 & 海光信息/寒武纪/中际旭创/新易盛 \\
    \addlinespace
    \textbf{功率半导体} & 92 & 35 & \textcolor{red}{$\bigstar\bigstar\bigstar$} & 35 & 48 & 32 & 25 & 4.2 & 新能源车+AI双驱动 & 海外厂商降价、产能释放 & 斯达半导/时代电气/扬杰科技/新洁能 \\
    \addlinespace
    \textbf{电力公用事业} & 78 & 18 & \textcolor{red}{$\bigstar\bigstar\bigstar$} & 12 & 22 & 16 & 32 & 2.1 & AI算力用电需求爆发 & 电价政策调整、煤价反弹 & 长江电力/国电南瑞/三峡能源/中国广核 \\
    \addlinespace
    \textbf{智能驾驶} & 82 & 42 & \textcolor{orange}{$\bigstar\bigstar$} & 28 & 35 & 45 & 45 & 3.5 & L3落地、城市NOA & 技术路线分歧 & 德赛西威/华阳集团/经纬恒润/伯特利 \\
    \addlinespace
    \textbf{人形机器人} & 72 & 30 & \textcolor{orange}{$\bigstar\bigstar$} & 55 & 85 & 58 & 28 & 1.8 & 特斯拉Optimus量产 & 量产节奏低于预期 & 绿的谐波/三花智控/拓普集团/鸣志电器 \\
    \addlinespace
    \textbf{光伏储能} & 68 & 15 & \textcolor{teal}{$\bigstar\bigstar$} & 18 & 25 & 22 & 12 & 5.5 & 储能需求爆发 & 价格战、贸易摩擦 & 阳光电源/锦浪科技/固德威/晶科能源 \\
    \addlinespace
    \textbf{创新药} & 75 & 48 & \textcolor{teal}{$\bigstar$} & 22 & 35 & 42 & 55 & 7.2 & GLP-1、ADC药物 & 临床失败、医保谈判 & 恒瑞医药/百济神州/信达生物/康方生物 \\
    \bottomrule
  \end{tabular}
  \vspace{0.3em}
  \caption{七大预期差板块核心指标横向对比}
  \label{tab:seven_surprise_sectors}
  \vspace{-0.5em}
  \begin{flushleft}
  \footnotesize\textit{数据来源：}Wind、各行业协会、本研究测算。数据截至2026年6月15日。
  \end{flushleft}
  \vspace{0.3em}
  \begin{tcolorbox}[colback=gray!5, colframe=gray!30, boxrule=0.5pt,
      left=6pt, right=6pt, top=4pt, bottom=4pt]
    \textbf{\scriptsize 预期差等级评定标准：}
    \scriptsize
    \begin{itemize}[leftmargin=*, itemsep=0pt, parsep=0pt]
      \item[\textcolor{red}{$\bigstar\bigstar\bigstar$ 强预期差}]：景气度评分$\geq$80且拥挤度评分$\leq$30，业绩与产业双验证，估值处于历史低位
      \item[\textcolor{orange}{$\bigstar\bigstar$ 中预期差}]：景气度评分70--80且拥挤度评分30--50，单一维度仍有分歧，存在预期修复空间
      \item[\textcolor{teal}{$\bigstar$ 弱预期差}]：景气度评分60--70或拥挤度评分50--60，边际改善趋势确立但弹性有限
    \end{itemize}
  \end{tcolorbox}
\end{table}
